\documentclass{article}
%packages
\usepackage{fullpage}
\usepackage{xspace}
\usepackage{graphicx,color}
\usepackage{hyperref}
\usepackage{graphicx}
\usepackage[dvipsnames]{xcolor}
\usepackage{amsmath, amsthm, amscd, amsfonts, amssymb}
\usepackage{graphicx}
\usepackage{stmaryrd}
\usepackage{quiver}
\usepackage{float}
\usepackage{mathtools}
\usepackage{mathrsfs}
\usepackage{multicol}
\usepackage{tikz-cd}
\usepackage[shortlabels]{enumitem}
\usepackage{stackrel}
\usepackage{cancel}
\usepackage[framemethod = TikZ]{mdframed}

%theorems
\newtheorem{theorem}{Teorema}[section]

\newtheorem{algorithm}{Algorithm}
\newtheorem{axiom}{Axioma}
\newtheorem{case}{Case}
\newtheorem{claim}{Afirmación}
\newtheorem{conclusion}{Conclusión}
\newtheorem{condition}{Condition}
\newtheorem{conjecture}[theorem]{Conjetura}
\newtheorem{corollary}[theorem]{Corolario}
\newtheorem{criterion}{Criterion}
\newtheorem{proposition}[theorem]{Proposición}
\newtheorem{lemma}[theorem]{Lema}

\theoremstyle{definition}
\newtheorem{definition}[theorem]{Definición}

\newtheorem{exercise}{Ejercicio}

\newtheorem*{notation}{Notación}
\newtheorem{problem}{Problema}[section]
\newtheorem*{obs}{Observación}
\newtheorem*{addendum}{Addendum}

\newtheorem*{tarea}{\wea{Tarea}}

\theoremstyle{remark}
\newtheorem{remark}[theorem]{Remark}
\newtheorem*{solution}{\textbf{Solución}}
\newtheorem{summary}{Summary}
\newtheorem{example}[theorem]{\textbf{Ejemplo}}

\numberwithin{equation}{section}

\renewcommand*{\proofname}{\textbf{Demostración}}

\surroundwithmdframed[
  topline=false,
  rightline=false,
  bottomline=false,
  leftmargin=\parindent,
  skipabove=\medskipamount,
  skipbelow=\medskipamount
]{proof}

\surroundwithmdframed[
  topline=false,
  rightline=false,
  bottomline=false,
  leftmargin=\parindent,
  skipabove=\medskipamount,
  skipbelow=\medskipamount
]{solution}

%commands
\newcommand{\A}{\mathbb{A}}
\newcommand{\B}{\mathcal{B}}
\newcommand{\C}{\mathbb{C}}
\newcommand{\D}{\mathbb{D}}
\newcommand{\F}{\mathbb{F}}
\newcommand{\M}{\mathcal{M}}
\newcommand{\N}{\mathbb{N}}
\renewcommand{\O}{\mathcal{O}}
\renewcommand{\P}{\mathcal{P}}
\newcommand{\Q}{\mathbb{Q}}
\newcommand{\R}{\mathbb{R}}
\renewcommand{\S}{\mathbb{S}}
\newcommand{\U}{\mathcal{U}}
\newcommand{\Z}{\mathbb{Z}}

\newcommand{\Acal}{\mathcal{A}}
\newcommand{\Bcal}{\mathcal{B}}
\newcommand{\Ccal}{\mathcal{C}}
\newcommand{\Dcal}{\mathcal{D}}
\newcommand{\Ecal}{\mathcal{E}}
\newcommand{\Fcal}{\mathcal{F}}
\newcommand{\Gcal}{\mathcal{G}}
\newcommand{\Hcal}{\mathcal{H}}
\newcommand{\Ical}{\mathcal{I}}
\newcommand{\Jcal}{\mathcal{J}}
\newcommand{\Kcal}{\mathcal{K}}
\newcommand{\Lcal}{\mathcal{L}}
\newcommand{\Mcal}{\mathcal{M}}
\newcommand{\Ncal}{\mathcal{N}}
\newcommand{\Qcal}{\mathcal{Q}}
\newcommand{\Rcal}{\mathcal{R}}
\newcommand{\Scal}{\mathcal{S}}
\newcommand{\Tcal}{\mathcal{T}}
\newcommand{\Ucal}{\mathcal{U}}
\newcommand{\Vcal}{\mathcal{V}}
\newcommand{\Wcal}{\mathcal{W}}
\newcommand{\Xcal}{\mathcal{X}}
\newcommand{\Ycal}{\mathcal{Y}}
\newcommand{\Zcal}{\mathcal{Z}}

\renewcommand{\a}{\alpha}
\renewcommand{\b}{\beta}
\renewcommand{\d}{\mathrm{d}}
\newcommand{\dd}{\delta} %% cambiar esto en los apuntes
\newcommand{\e}{\varepsilon}
\renewcommand{\t}{\tau}
\newcommand{\lam}{\lambda}

\newcommand{\vphi}{\varphi}
\newcommand{\oo}{\infty}
\newcommand{\isom}{\cong}
\newcommand{\longto}{\longrightarrow}
\newcommand{\embbed}{\hookrightarrow}

\renewcommand{\bar}{\overline}
\newcommand{\sgn}{\mathrm{sgn}}
\newcommand{\id}{\mathrm{id}}
\newcommand{\im}{\mathrm{im}\,}
\newcommand{\dist}{\mathrm{dist}}

\newcommand{\normal}{\unlhd}

\renewcommand{\tilde}{\widetilde}
\renewcommand{\hat}{\widehat}

\newcommand{\abs}[1]{\left| #1 \right|}
\newcommand{\norm}[1]{\left\lVert #1 \right\rVert}
\newcommand{\br}[1]{\left( #1 \right)}
\newcommand{\set}[1]{\left\{ #1 \right\}}
\newcommand{\cor}[1]{\left[ #1 \right]}
\newcommand{\gen}[1]{\left\langle #1 \right\rangle}
\newcommand{\matr}[1]{\begin{pmatrix}#1\end{pmatrix}}

\newcommand{\znz}[1]{\frac{\Z}{#1 \Z}}
\newcommand{\zmz}[1]{\Z/ #1 \Z}

% Por si me equivoco al escribir
\newcommand{\rac}[2]{\frac{#1}{#2}}
\newcommand{\olon}{\colon}
\newcommand{\ubseteq}{\subseteq}
\newcommand{\ubset}{\subset}
\newcommand{\irc}{\circ}
\renewcommand{\o}{\to}

\newcommand{\ds}{\displaystyle}

\renewcommand*\contentsname{Contenidos}

\newcommand{\wea}[1]{\textcolor{BlueViolet}{\textbf{\emph{#1}}}}
\newcommand{\fecha}[1]{\begin{flushright} \underline{#1} \end{flushright}}

\newcommand{\dem}[1]{\textbf{Demostración (#1)}}

%title
\title{MPG3200 -- Análisis I \\ Ayudantías}
\author{Felipe Inostroza \\ Prof. Mariel Sáez \\ Ayud. Fernando Castillo}
\date{Primer Semestre 2026}

\begin{document}

\maketitle
\tableofcontents

\newpage
\section{Ayudantía 1. Lunes 16 de Marzo}
\subsection{Problema 1}

\begin{enumerate}
\item Sea $X$ un conjunto y $\t_f$ la colección de subconjuntos $U$ de $X$ tal que $U^c$ es finita o todo $X$. Muestre que $(X,\t_f)$ es un espacio topológico. Llamamos a $\t_f$ la topología cofinita cofinita en $X$.
\begin{proof}
    Notemos que $\varnothing, X \in \t_f$. Consideremos $\{U_\a\}_\a \subset \t_f$, entonces tenemos que $U_\a^c$ son finitos $\forall \a$, y luego $\ds \bigcap_\a U_\a^c = \br{\bigcup_\a U_\a}^c$ es finito. Luego $\bigcup_\a U_\a \in \t_f$.
    Sean $U_1,\dots,U_n \in \t_f$, tenemos que $U_i^c$ son finitos, y por lo tanto $\bigcup_{i \leq n} U_i^c = \br{\bigcap_{i\leq n} U_i}^c$ es finito, y entonces $U_1 \cap \dots \cap U_n \in \t_f$.
\end{proof}

\item Muestre que si $X$ no es numerable, entonces $(X,\t_f)$ no es primero contable.
\begin{proof}
    Supongamos que $(X,\t_f)$ es primero contable. Sea $x \in X$ y $\{\B_i\}_{i\in\N}$ base numerable en $x$.
    Consideremos el conjunto
    \[F = \{x\} \cup \bigcup_{n>0} B_n^c.\]
    Notemos que $F$ es numerable, por lo que $X \setminus F \neq \varnothing$.
    Sea $y \in X \setminus F$, entonces $\{y\}^c$ es vecindad de $x$. Luego existe $n_0 \in \Z_{>0}$ tal que $x \in B_{n_0} \subseteq \{y\}^c$, y entonces $y \in B_{n_0}^c \subseteq F$. $\to\gets$

    Por lo tanto $(X,\t_f)$ no es primero contable.
\end{proof}

\item Denotando $\R_f$ como $\R$ con topología cofinita, evalúe a qué puntos converge $x_n = \dfrac 1n$.
\begin{proof}
    Mostremos que $x_n$ converge a todo $\R$.
    Sea $x \in \R$, y sea $U$ vecindad de $x$. Escribimos $$U^c = \{y_1,\dots,y_m\}.$$
    Luego existe $n_0 \in \Z_{>0}$ tal que $x_n \in U$ para todo $n \geq n_0$, ya que de caso contrario $U^c$ sería infinito.
\end{proof}
\end{enumerate}

\newpage
\subsection{Problema 2}
\begin{enumerate}
  \item En $\R^n$, considere $\B$ la familia de conjuntos $\{x \in \R^n \colon p(x) \neq 0\}$, donde $p(x)$ es un polinomio de $n$ variables. MMuestre que $\B$ define una base para una topología. Llamamos a la topología generada la topología de Zariski en $\R^n$.
  \begin{proof}
    \
    \begin{enumerate}[(i)]
      \item Denotamos $C(p) = \{x \in \R^n \colon p(x) \neq 0\}$. Notemos que
      \[\bigcup_{B \in \B} B = \bigcup_{p \in \R[x_1,\dots,x_n]} C(p) \supseteq C(1) = \{x \in \R^n \colon 1 \neq 0\} = \R^n.\]
      \item Sean $C(p), C(q) \in \B$ y $x \in C(p) \cap C(q)$. Tenemos que $p(x) \neq 0 \neq q(x)$. Luego $p(x) q(x) \neq 0$. Entonces $x \in C(pq)$. Por lo mismo, $C(pq) \subseteq C(p) \cap C(q)$.
    \end{enumerate}
  \end{proof}

  \item Muestre que $\R^n$ en topología de Zariski es Tychonoff pero no Hausdorff.
  \begin{proof}
    Supongamos que es Hausdorff. Luego $\R \times \{0\}^{n-1} \subseteq \R^n$ debiese ser Hausdorff. La base para $\R$ es formada por conjuntos de la forma
    \[\{(x,0,\dots,0) \in \R^n \colon p(x,0,\dots,0) \neq 0\}.\]
    Es decir, $\R$ como subespacio coincide con $\R$ en topología de Zariski. En $\R$
    \[C(p) = \{x \in \R \colon p(x) \neq 0\} = \{z_!,\dots,z_n\}^c\]
    esto por el Teorema Fundamental del Álgebra. De aquí, $\R$ en Zariski coincide con $\R_f$.
    En espacios Hausdorff, los límites son únicos, pero por el Problema 1, $\R_f$ no cumple esto. $\to\gets$
  \end{proof}
\end{enumerate}

\newpage
\subsection{Problema 3}
Muestre que $X$ es un espacio topológico Hausdorff si y solo si $$\Delta = \{(x,x) \colon x \in X\} \subseteq X \times X$$ es cerrado.
\begin{proof}
  \
  \begin{itemize}
    \item [\boxed{\implies}]
    Supongamos que $X$ es Hausdorff. Veamos que $\Delta^c$ es abierto.
    Sea $(x,y) \in \Delta^c$, entonces $x \neq y$. Como $X$ es Hausdorff, existen abiertos $U,V$ disjuntos tales que
    \[x \in U, \quad y \in V.\]
    Como $U,V$ son disjuntos, tenemos que $(t_1,t_2) \in U \times V \implies t_1 \neq t_2$, y entonces $(U \times V) \cap \Delta = \varnothing$. Por lo tanto tenemos que $$(x,y) \in U \times V \subseteq \Delta^c.$$
    Luego por cada punto en $\Delta^c$ existe una vecindad dentro de $\Delta^c$, y entonces $\Delta^c$ es abierto. Luego $\Delta$ es cerrado.

    \item [\boxed{\impliedby}]
    Supongamos que $\Delta$ es cerrado. Sen $x \neq y$ en $X$, entonces $(x,y) \in \Delta^c$. Sean $U,V$ abiertos con $$(x,y) \in U \times V \subseteq \Delta^c.$$
    Luego $x \in U, \ y \in V$. Dado que $U \times V \subseteq \Delta^c$, $U$ y $V$ deben ser disjuntos. Por lo tanto $X$ es Hausdorff.
  \end{itemize}
\end{proof}

\newpage
\subsection{Problema 4}
La recta de Sorgenfrey $\R_l$ se define como $\R$ con la topología generada por la base $[a,b)$ para $a<b$.
\begin{enumerate}
  \item Muestre que $\R_l$ es primero contable, separable pero no segundo contable.
  \item Concluya que $\R_l$ no es metrizable.
\end{enumerate}

\newpage
\subsection{Problema 5}
Sea $K = \{1/n \colon n \in \N\}$. La $K$--topología en $\R$ se obtiene a partir de la base dada por intervalos abiertos $(a,b)$ y conjuntos de la forma $(a,b) \setminus K$.
Muestre que $\R$ con la $K$--topología es Hausdorff pero no regular.
\begin{proof}
  \begin{itemize}
    \item \wea{Hausdorff:}
    Sean $x,y \in \R$ distintos. Como $\R$ con la topología euclideana es Hausdorff, existen $U,V$ disjuntos tales que $x \in U$ e $y \in V$. Como la $K$--topología contiene la Euclideana, $U,V$ son abiertos en $K$--topología. Por lo tanto es Hausdorff.

    \item \wea{No regular:}
    Tenemos que $\R \setminus K$ es abierto, y entonces $K$ es cerrado. También tenemos que $0 \notin K$. Supongamos que la $K$--topología es regular, entonces existen $U,V$ abiertos disjuntos tales que $0 \in U$ y $K \subseteq V$. Usando la base, tenemos que existen $a,b \in \R$ tales que
    \[0 \in (a,b) \subseteq U \quad\lor\quad 0 \in (a,b) \setminus K \subseteq U.\]
    No me acuerdo como se concluía pero tenemos que $0 \in (a,b) \subseteq U$.

    Para $K$, tomamos $1/m \in (a,b)$, luego existen $c,d \in \R$ tales que
    \[\frac 1m \in (c,d) \subseteq V \quad\lor\quad \frac 1m \in (c,d) \setminus K \subseteq V.\]
    Notemos que $1/m \in K$, por lo que $1/m \notin (c,d) \setminus K$, y entonces $1/m \in (c,d) \subseteq V$. Luego podemos pedir
    $$(c,d) \subseteq (a,b) \quad\lor\quad (c,d) \subseteq V$$
    Si $z \in (c,d)$ es irracional, entonces $z \in (a,b) \setminus K \subseteq U$. Luego $U \cap V \neq \varnothing$, y entonces la $K$--topología no es regular.
  \end{itemize}
\end{proof}

\newpage
\section{Ayudantia 2. Lunes 23 de Marzo}

\subsection{Problema 1}

La recta de Sorgenfrey $\R_f$ es $\R$ con la base $[a,b)$ con $a<b$.
\begin{enumerate}
  \item Muestre que $\R_f$ no es segundo numerable.
  \begin{proof}
    Sea $\B$ una base de $\R_f$. Sea $x \in \R$, tenemos que $[x,x+1)$ es una vecindad abierta de $x$.
    Luego existe $B_x \in \B$ tal que $x \in B_x \subseteq [x,x+1)$.
    Si $y \neq x$, hacemos lo mismo:
    \[y \in B_y \subseteq [y,y+1).\]
    Vemos que $B_x \neq B_y$, pues $\inf{B_x} \geq \inf{[x,x+1)} = x$, y como $x \in B_x$, entonces $\inf{B_x} = x$. Similarmente, $\inf{B_y} = y$, y por lo tanto $\inf{B_x} \neq \inf{B_y}$. Luego $B_x \neq B_y$.

    Entonces tenemos una inyección $\R \to \B$ dada por $x \mapsto B_x$, y entonces $\B$ no puede ser numerable.
  \end{proof}

  \item Muestre que $\R_f$ es Lindelöf.
  \begin{proof}
    Mostraremos que todo subrimiento por elementos de la base tiene subcubrimiento numerable.

    Sea $\Acal = \{[a_\a,b_\a)\}_\a$ cubrimiento, y definimos
    \[C = \bigcup_\a (a_\a,b_\a).\]
    Veamos que $\R \setminus C = \{a_\a\}_\a$ es numerable. Sea $x \in \R$, como $\Acal$ es cubrimiento, entonces $\exists \b$ tal que $$x = a_\b \in [a_\b,b_\b)$$.
    Sea $q_x \in \Q \cap [a_\b,b_\b) = \Q \cap [x,b_\b)$. Si $x \neq y$, entonces podemos asumir $x < y$, y entonces $q_x < q_y$, ya que en caso contrario tendriamos $x < y < q_y \leq q_x$ por lo que $y \in (x,q_x) \subseteq (a_\b,b_\b) \subseteq C$, pero $y \notin C$.

    Tenemos entonces una inyección $\R \setminus C \to \Q$ dada por $x \mapsto q_x$, y entonces $\R \setminus C$ es numerable.

    Por otro lado, viendo $C$ con topología Euclideana, al ser segundo numerable, admite subcubrimiento numerable $\{(a_{\a_n},b_{\a_n})\}_{n\in\N}$. Luego $\{[a_{\a_n},b_{\a_n})\}_{n\in\N} \cup \R\setminus C$ es numerable.
  \end{proof}

  \item Muestre que $\R_f \times \R_f$ no es Lindelöf.
  \begin{proof}
    Consideremos la recta $L = \{(x,-x) \colon x \in \R\}$. Tomamos
    \[\Acal = \{L^c\} \cup \{[x,x+1) \times [-x,-x+1) \colon x \in \R\}.\]
    De ser Lindelöf, cubriríamos $L$ con numerables elementos de $\Acal$, pero cada elemento de $\Acal$ contiene a lo más un elemento de $L$. Pero $\abs L = \abs \R$. $\to\gets$
  \end{proof}
\end{enumerate}

\newpage
\subsection{Problema 2}

\begin{enumerate}
  \item Si $N \subseteq X \times Y$ es abierto con $Y$ compacto y $\{x_0\} \times Y \subseteq N$, existe $W \subseteq X$ vecindad de $x_0$ tal que $$\{x_0\} \times Y \subseteq W \times Y \subseteq N.$$
  \begin{proof}
    Cubrimos $\{x_0\} \times Y$ con elementos de la base $U \times V \subseteq N$. Notemos que $\{x_0\} \times Y$ es homeomorfo a $Y$, el cual es compacto, y entonces existe un cubrimiento finito $U_1 \times V_1, \dots, U_n \times V_n$.
    Sea $W = U_1 \cap \dots \cap U_n$.
    
    Veamos que $\{U_i \times V_i\}_{i=1}^n$ cubre a $W \times Y$.
    Sea $(x,y) \in W \times Y$. Consideramos $(x_0,y) \in \{x_0\} \times Y$. Sea $U_i \times V_i$ con $(x_0,y) \in U_i \times V_i$, entonces $y \in V_i$. Notemos que
    \[x \in W = \bigcap_{j=1}^n U_j \subseteq U_i\]
    y entonces $(x,y) \in U_i \times V_i$. Por lo tanto
    \[\{x_0\} \times Y \subseteq W \times Y \subseteq \bigcup_{i=1}^n U_i \times V_i \subseteq N.\]
  \end{proof}
  
  \item Muestre que si $Y$ es compacto, entonces la proyección $\pi_1 \colon X \times Y \to X$ es cerrada.
  \begin{proof}
    Sea $F \subseteq X \times Y$ cerrado. Queremos mostrar que $\pi_1(F)$ es cerrado, lo cual pasa si y solo si $\pi_1(F)^c$ es abierto. Sea $x_0 \in \pi_1(F)^c \subseteq X$, entonces $\{x_0\} \times Y \subseteq F^c \subseteq X \times Y$.
    Luego por la parte anterior, existe $W \subseteq X$ vecindad de $x_0$ tal que $$\{x_0\} \times Y \subseteq W \times Y \subseteq F^c.$$
    Luego
    \[x_0 \in W = \pi_1(W \times Y) \subseteq \pi_1(F^c).\]
    Faltó concluir, pero el Fer mandará la solución después.
  \end{proof}
\end{enumerate}

\newpage
\section{Ayudantía 3. Lunes 30 de Marzo}

\subsection{Problema 1}
Sea $X$ un espacio conexo normal. Muestre que $X$ es un singleton o bien no numerable.
\begin{proof}
  Si $X = \{*\}$ estamos listos. Espléndido.

  Si $\abs{X}>1$, sean $a,b \in X$ con $a \neq b$. Por Urysohn, existe $f \colon X \to [0,1]$ continua tal que $f(a) = 0$ y $f(b) = 1$. Como $X$ es conexo, entonces $f(X) \subseteq [0,1]$ también es conexo, y entonces $f(X)$ es un intervalo. Pero $0,1 \in f(X)$, y entonces $f(X) = [0,1]$. Luego, como $f$ es sobreyectiva:
  \[\abs\R = \abs{[0,1]} \leq \abs X.\]
\end{proof}

\newpage
\subsection{Problema 2}
$X$ es conexo por caminos si para todo $p,q \in X$, existe $\gamma \colon [0,1] \to \R$ continua tal que $\gamma(0) = p$ y $\gamma(1) = q$.
\begin{enumerate}
\item Muestre que si $U \subseteq \R^n$ es un abierto conexo, entonces es conexo por caminos.
\begin{proof}
  Sea $p \in U$, y sea
  \[V = \{q \in U \mid \exists \gamma \colon [0,1] \to U, \quad \gamma(0) = p, \quad \gamma(1) = q\}.\]
  Veamos que $V$ es abierto cerrado no vacío, así concluiremos por conexidad que $V = U$.

  \begin{itemize}
    \item Notemos que $p \in V$, por lo que $V \neq \varnothing$.
    \item \wea{Abierto: }
    Sea $q \in V \subseteq U$. Sea $B$ tal que $q \in B \subseteq U$.
    Queremos $q \in B \subseteq V$. Pero las bolas son convexas, por lo que $q \in B \subseteq V$, y entonces $V$ es abierto.
    \item \wea{Cerrado: }
    Tomar $q \in \partial V$, no me acuerdo como seguía fkgjfgk
  \end{itemize}
\end{proof}

\item Muestre que si $A \subseteq \R^2$ es numerable, entonces $\R^2 \setminus A$ es conexo por caminos.
\begin{proof}
  Sean $p,q \in \R^2 \setminus A$. Primero notemos que existe una recta $L_p$ tal que $p \in L_p \subseteq \R^2 \setminus A$, esto porque existe una cantidad no numerable de rectas que pasan por $p$, y como $A$ es numerable, solo hay una cantidad numerable de rectas que pasan por $p$ y que pasan por algún punto de $A$. Entonces existe la recta $L_p$.

  Similarmente, podemos encontrar una recta $L_q$ tal que $q \in L_q \subseteq \R^2 \setminus A$ y que además no sea paralela a $L_p$ (como tenemos no numerables rectas contenidas en $\R^2 \setminus A$, solo una de esas podría ser paralelas, y tenemos infinitas que no lo son). Luego $L_p \cap L_q \neq \varnothing$, y luego tenemos el camino dado por recorrer desde $p$ hasta la intersección, y después ir hasta $q$.
\end{proof}
\end{enumerate}

\newpage
\subsection{Problema 3}
Sea $X$ un espacio topológico normal, $F \subseteq X$ un cerrado y $g \colon F \to \R$ continua. Siga el siguiente procedimiento para mostrar que existe una extensión de $g$ a todo $X$. (¿Por qué no es consecuencia inmediata de Tietze?)
\begin{enumerate}
\item Reduzca el caso a $f \colon F \to (-1,1)$.
\begin{proof}
  Tenemos un homeomorfismo $\psi \colon \R \to (-1,1)$, y entonces tomamos $f = \psi \circ g \colon F \to \R \to (-1,1)$ continua.
\end{proof}

\item Extienda $f$ a una función $h$ en todo $X$.
\begin{proof}
  Como $(-1,1) \subseteq [-1,1]$, tenemos $f \colon F \to [-1,1]$, y por Tietze existe $h \colon X \to [-1,1]$ tal que $h\vert_F = f$.
\end{proof}

\item Consiga una función $\vphi \colon X \o [-1,1]$ tal que $\vphi = 1$ en $F$ y $\vphi = 0$ en $h^{-1}(\{-1,1\})$.
\begin{proof}
  Por Urysohn, existe $\vphi \colon X \to [0,1]$ tal que $\vphi(F) = \{1\}$ y $\vphi(h^{-1}(\{-1,1\})) = \{0\}$.
  Además, $F \cap h^{-1}(\{-1,1\}) = \varnothing$, pues $h(x) = f(x) \in (-1,1)$ para todo $x \in F$.
\end{proof}

\item Considere $\bar g = \vphi \cdot h$.
\begin{proof}
  Considere Considere $\bar g = \varphi \cdot h \colon X \to (-1,1)$. Veamos que extiende a $g$: Para $x \in F$ tenemos
  \[\bar g(x) = \vphi(x) \cdot h(x) = 1 \cdot g(x) = g(x).\]
  Por otro lado, si $x \in X$, tenemos
  \begin{align*}
    &\vphi(x) h(x) = 1 \\
    \iff &\vphi(x) = h(x) = 1 \\
    \implies &x \in h^{-1}(\{1\}) \\
    \implies &\vphi(x) = 0. \quad \to\gets
  \end{align*}
  Por lo tanto $\bar g \colon X \to [-1,1)$. Mismo arumento para $\bar g \colon X \to (-1,1)$.
\end{proof}
\end{enumerate}

\newpage
\subsection{Problema 4}
Dado $X$ un espacio topológico, un \wea{retracto} es un subespacio $Y \subseteq X$ tal que existe una función continua $r \colon X \to Y$ con $r(y) = y$ para todo $y \in Y$.

Sea $X$ un espacio normal e $Y \subseteq X$ un cerrado tal que $Y \simeq \R$. Muestre que $Y$ es un retracto de $X$.

\begin{proof}
  Consideremos $\id_Y \colon Y \to Y$. Sea $f \colon Y \xrightarrow\sim \R$, luego $f \circ \id_Y \colon Y \to \R$. Por el problema 3, existe $\bar{f \circ \id_Y} \colon X \to \R$ extensión continua.
  Consideramos $r = f^{-1} \circ \bar{f \circ \id_Y} \colon X \to Y$.

  Para $y \in Y$ tenemos
  \begin{align*}
    r(y) &= f^{-1}(\bar{f \circ \id_Y}(y)) \\
    &= f^{-1}(f \circ \id_Y(y)) \\
    &= f^{-1} \circ f (y) = y.
  \end{align*}
  Por lo tanto $r$ es retracto.
\end{proof}

\newpage
\subsection{Problema 5}
El plano de Moore $$\Gamma = \{(x,y) \in \R^2 \colon y \geq 0\}$$ se le equipa la topología dada por la siguiente base local: Si $(x,y)$ con $y > 0$, los elementos de la base son $B_r((x,y)) \subseteq \Gamma$. Si $(x,0) \in \Gamma$, los elementos de la base son $\{(x,0)\} \cup B_r((x,r))$, discos abiertos tangentes a $(x,0)$ con $(x,0)$.

Muestre que el plano de Moore es regular no normal.
\begin{proof}
  Consideremos $D = (\Q_{>0} \times \Q_{>0}) \cap \{y>0\}$. Este conjunto es denso en $\Gamma$. Luego las funciones continuas $C(\Gamma)$ quedan determinadas por $f \colon D \to \R$. Entonces a lo más hay $\abs{\R}^{\abs D} = \abs{\R}^{\abs\N}$ funciones continuas reales.

  $L = \{(x,0) \inf\Gamma \colon x \in \R\}$ es el complemento del seiplano superior, por lo que es cerrado, y también con la topología subespacio es discreto. Luego las funciones continuas $g \colon L \to \R$ se extienden a $\bar g \colon \Gamma \to \R$, y entonces hay por lo menos $\abs\R^{\abs\R}$ funciones continuas. $\to\gets$
\end{proof}

\newpage
\section{Ayudantía 4}
\subsection{Problema 1}

Sea $f \colon [a,b] \to \R$ continua. Si $\int_a^b f(x) x^n \d x = 0$ para todo $n \in \N$, entonces $f$ es idénticamente nula.
\begin{proof}
  Por Stone--Weierstrass, $p_k \to f$ para polinomios $p_k$.

  Si $q(x) = \sum_{i=1}^N a_i x^i$ es un polinomio, entonces
  \[\int_a^b f(x) q(x) \d x = \sum_{i=1}^N a_i \int_a^b f(x) \cdot x^i \d x = 0\]
  \[\int_a^b f^2(x) \d x
  = \int_a^b f(x) f(x) \dashcolon= \lim_{k\to\oo} \int_a^b f(x) p_k(x) \d x = 0\]
  Luego $f(x) = 0$ para todo $x \in [a,b]$.
\end{proof}

\newpage
\subsection{Problema 2}

Sean $M,N \subseteq \R$ compactos Hausdorff. Toda función continua $f \colon M \times N \to \R$ se puede aproximar uniformemente por sumas de funciones del tipo $(x,y) \mapsto u(x)v(y)$ con $u \colon M \to \R$ y $v \colon N \to \R$ continuas.

\begin{proof}
  Definimos
  \[A = \set{\sum_{k=1}^d u_k(x) v_k(y) \colon u_k \in C(M), \ v_k \in C(N)}.\]
  Notemos que $A$ es un $\R$--espacio vectorial. Luego
  \[\br{\sum_{k=1}^{d_1} u_k(x)v_k(y)}\br{\sum_{l=1}^{d_2} \tilde u_l(x) \tilde v_l(y)} \in A\]
  No voy a tipear ese desarrollo.
  Luego $A$ es una álgebra de funciones. Notemos también que si $c \in \R$, entonces $c \in C(M)$ y $1 \in C(N)$ nos da que $c \in A$.

  Separa puntos:
  Sean $(x,y) \neq (w,z)$. Supongamos que $x \neq w$. Como $X$ es compacto Hausdorff, $X$ es normal. Luego por Urysohn, $\exists f \colon M \to [0,1]$ continua tal que $f(x) = 1$ y $f(w) = 0$. Basta con $u(x) = f(x)$ y $v(y) = 1$. la función $u(x)v(y) = f(x)$ separa a $(x,y)$ de $(w,z)$. Por lo tanto $A$ separa puntos.

  Finalmente, por Stone--Weierstrass tenemos que $A$ es denso en $C(X)$. Por lo tanto $f \colon M \times N \to \R$ se pueden aproximar unirofmemente por funciones de $A$.
\end{proof}

\newpage
\subsection{Problema 3}
Sea $X$ un espacio  métrico compacto y $\Fcal \subseteq C(X)$.
\begin{enumerate}
\item Muestre que $\Fcal$ es equicontinua si y solo si $\bar\Fcal$ es equicontinua.
\begin{proof}
  Si $\bar\Fcal$ es equicontinua, como $\Fcal \subseteq \bar\Fcal$, tenemos que $\Fcal$ es equicontinua.
  Si $\Fcal$ es equicontinua:

  Sea $f \in \bar\Fcal$, y sea $\{f_k\}_k$ tal que $f_k \xrightarrow{d_\oo} f$. Sea $\e>0$, entonces existe $\dd>0$ tal que
  \begin{align*}
    d(x,y) < \dd &\implies \abs{f_k(x) - f_k(y)} < \frac\e2 \\
    &\implies \lim_{k\to\oo} \abs{f_k(x) - f_k(y)} \leq \frac\e2 \\
    &\implies \abs{f(x) - f(y)} < \e.
  \end{align*}
  Como $f \in \bar\Fcal$ fue arbitrario, $\bar\Fcal$ es equicontinua.
\end{proof}

\item $\Fcal \subseteq C(X)$ es compacto si y solo si $\Fcal$ es una familia cerrada, uniformemente acotada y equicontinua.
\begin{proof}
  Visto en clases.
\end{proof}

\item Muestre que un subconjunto de $C(X)$ tiene clausura compacta si y solo si es equicontinua y uniformemente acotada.
\begin{proof}
  \begin{itemize}
    \item \boxed{\impliedby} Arzela--Ascoli.
    \item \boxed{\implies} $\Fcal$ tiene clausura compacta, entonces $\bar\Fcal$ es compacto. Luego $\bar\Fcal$ es equicontinua uniformemente acotara $\implies \Fcal$ es equicontinua uniformemente acotada.
  \end{itemize}
\end{proof}
\end{enumerate}


\newpage
\subsection{Problema 4}
Sea $K \in C([0,1] \times [0,1])$. Para $f \in C([0,1])$, sea $Tf(x) = \int_0^1 K(x,y)f(y) \d y$. Entonces, $Tf \in C([0,1])$ y $\Fcal = \{Tf \colon \norm{f}_\oo \leq 1\}$ tiene clausura compacta.

\begin{proof}
  Notemos que
  \begin{align*}
    &\abs{\int_0^1 (K(x+h,y) - K(x,y)) f(y) \d y} \\
    &\leq \int_0^1 \abs{K(x+h,y) - K(x,y)} \cdot \abs{f(y)} \d y \\
    &\leq \norm{f}_\oo \int_0^1 \abs{K(x+h,y) - K(x,y)} \d y \\
    &< \norm{f} \cdot 1 = \e
  \end{align*}
  Por continuidad uniforme de $K$. Luego $Tf \in \C([0,1])$.

  Veamos que $\Fcal$ tiene clausura compacta.
  Sea $\ds M = \max_{(x,y) \in [0,1]^2} K(x,y)$
  \begin{align*}
    \abs{Tf(x)} &= \abs{\int_0^1 K(x,y) f(y) \d y} \\
    &\leq \int_0^1 \abs{K(x,y) f(y)} \d y \\
    &\leq \int_0^1 M \norm{f}_\oo \d y \\
    &= M \norm{f}_\oo \leq M.
  \end{align*}
  Por lo tanto $\norm{Tf}_\oo \leq M$. Luego es uniformemente acotada.

  \begin{align*}
    \abs{Tf(x) - Tf(z)} 
    &= \abs{\int_0^1 (K(x,y) - K(z,y)) f(y)} \d y \\
    &\leq \norm{f}_\oo \int_0^1 \abs{K(x,y) - K(z,y)} \d y \\
    &= \int_0^1 \abs{K(x,y) - K(z,y)} \d y
  \end{align*}
  $\forall \e>0, \exists \dd>0$ tal que
  \[\abs{x-z} < \dd \implies \abs{K(x,y) - K(z,y)} < \e\]
  y en tal caso
  \[\abs{Tf(x) - Tf(z)} \leq \int_0^1 \frac\e2 \d y = \frac\e2 < \e\]
  Luego como $f \in \Fcal$ arbitrario, tenemos que $\Fcal$ es equicontinua.
\end{proof}

\newpage
\subsection{Problema 5}

Sea $\{f_n\}_n \subseteq C([a,b])$.

\begin{enumerate}
\item Si $\{f_n\}_n$ es equicontinua, ¿necesariamente debe tener una subsucesión uniformemente convergente?
\begin{proof}
  Sean $f_n \colon [0,1] \to \R$ con $x \mapsto x^n$. Esta sucesión de funciones es uniformemente acotada, pues $\abs{f_n(x)} \leq 1$ para todo $x \in [0,1]$.
  Notemos que $f_n \to f$ donde
  \[f(x) = \begin{cases}
    0, &x \in [0,1) \\
    1, &x \in \{1\}
  \end{cases}\]
  la cual no es continua, así que no puede ser uniforme la convergencia.
\end{proof}

\item Si $\{f_n\}_n$ es uniformemente acotada, ¿necesariamente debe tener una subsucesión uniformemente convergente?
\begin{proof}
  No, Consideremos $f_n(x) = n$.
\end{proof}
\end{enumerate}

\newpage
\section{Ayudantía 5. Miércoles 22 de Abril}

\subsection{Problema 0}

Sea $X$ completo. No existe $A \subseteq X$ tal que $A$ y $A^c$ sean $G_\delta$ densos.

\begin{proof}
  Por contradicción, supongamos que existe $A \subseteq X$ tal que $A$ y $A^c$ sean $G_\delta$ densos. Entonces existe $\{U_n\}_n \subseteq \t$ tal que 
  \[ A = \bigcap_{n \in \N} U_n. \]
  Luego para todo $n$ tenemos 
  \[ A \subseteq U_n \implies \bar A \subseteq \bar U_n \implies X = \bar U_n \]
  y entonces $U_n$ es denso para todo $n \in \N$. De la misma manera, existe $\{V_n\}_n \subseteq \t$ tal que 
  \[ A^c = \bigcap_{n \in \N} V_n  \]
  y por el mismo argumento $V_n$ también es denso $\forall n \in \N$. Luego $\{U_n\}_{n\in\N} \cup \{V_n\}_{n\in\N}$ es una colección de abiertos densos. Luego 
  \[ \varnothing= A \cap A^c = \left( \bigcap_{n\in\N} U_n \right) \cap \left( \bigcap_{n \in \N} V_n  \right) \]
  y entonces $\varnothing$ sería denso. \(\to\gets\)
\end{proof}

\newpage
\subsection{Problema 1}

Sea $X$ un espacio métrico y $f \colon X \to \R$. Muestre que los puntos en que $f$ es continua son intersecciones numerables de abiertos. Concluya que no hay funciones $f \colon \R \to \R$ continuas solamente en los racionales.

\begin{proof}
  Sea $\Ccal \subseteq X$ el conjunto de los puntos donde $f$ es continua. Queremos ver que $\Ccal$ sea $G_\delta$. 
  \begin{align*}
      \Ccal &= \{x_0 \in X \colon \forall \e > 0, \exists \delta > 0, \forall y \in B_\delta(x_0) \implies \abs{f(y) - f(x_0)} < \e\}
      % \\ &= \bigcap_{n\in\N} \set{x_0 \in X \colon \exists \delta>0, \forall y \in B_\delta(x_0), \abs{f(y) - f(x_0)} < 1/n}
  \end{align*}
  Tomamos 
  \[ A_n = \set{x_0 \in X \colon \exists \delta>0, \forall y,z \in B_\delta(x_0), \abs{f(y) - f(z)} < \frac1n} \]
  Queremos mostrar 
  \[ \Ccal = \bigcap_{n\in\N} A_n. \]
  Notemos que $A_n$ es abierto (demostrado por dibujo). Sea $n \in \N$, tomando $z = x_0$, tenemos $\forall n \in \N$, $\exists \delta_n>0$, $\forall y \in B_{\delta_n}(x_0)$, 
  \[ \abs{f(y) - f(x_0)} < \frac1n. \]

  Sea $x_0 \in \Ccal$, queremos que $x_0 \in A_n$ para todo $n$. Tenemos que, con $\e = 1/2n$, $\exists \delta>0$, $\forall y \in B_\delta(x_0)$, 
  \[ \abs{f(y) - f(x_0)} < \frac{1}{2n}. \]
  Luego, si $z \in B_\delta(x_0)$, tenemos que 
  \[ \abs{f(y) - f(z)} \leq \abs{f(y) - f(x_0)} + \abs{f(x_0) - f(z)} < \frac{1}{2n} + \frac{1}{2n} = \frac{1}{n}. \]
  Por lo tanto $\Ccal \subseteq \bigcap_n A_n \subseteq \Ccal \implies \Ccal = \bigcap_n A_n$. Luego $\Ccal$ es $G_\delta$.

  Ahora, supongamos que existe $f \colon \R \to \R$ tal que $\Ccal = \Q$. Luego $\Q$ es un $G_\delta$, y además es denso. Por otro lado, $\R \setminus \Q$ es denso. También, $\Q = \bigcup_{q\in\Q} \{q\}$, por lo que es $F_\sigma$, y entonces $\R \setminus \Q$ es $G_\delta$. Es decir, $\Q$ y $\R \setminus \Q$ son $G_\delta$ densos, y por el problema anterior esto no puede pasar.
\end{proof}

\newpage
\subsection{Problema 2}
Sea $\{x_n\}_n \subseteq \R$ sucesión. Defina $\mu$ en $(\R, \B(\bar\R))$ por $\mu = \sum_n \delta_{x_n}$.

\begin{enumerate}
\item Muestre que $\mu$ asigna finitos valores en intervalos acotados de $\R$ si y solo si $\lim_{n\to\oo} \abs{x_n} = \oo$.

\item ¿Para qué sucesiones $\{x_n\}_n$ es $\mu$ una medida $\sigma$--finita?
\end{enumerate}

\newpage
\subsection{Problema 3}
Sea $(X,\M,\mu)$ un espacio de medida.
\begin{enumerate}
  \item Muestre que si $E$ y $F$ son medibles, entonces $\mu(E \triangle F) = 0$ implica que $\mu(E) = \mu(F)$.
  \item Muestre que si $\mu$ es completa ($A \subseteq E$ con $E \in \M$ y $\mu(E) = 0$, entonces $A \in \M$), $E \in \M$ y $\mu(E \triangle F) = 0$, entonces $F \in \M$.
  \item Defina en $\M$ la relación $E \sim F$ si y solo si $\mu(E \triangle F) = 0$. Muestre que $\sim$ es relación de equivalencia.
  \item Muestre que $\rho(E,F) = \mu(E \triangle F)$ cumple $\rho(E,F) \leq \rho(E,G) + \rho(G,F)$. (Con esto, acaba de equipar una métrica al espacio $\M/{\sim}$).
\end{enumerate}
\begin{proof}
  \ \begin{enumerate}
    \item 
    \begin{align*}
        \mu(E) &= \mu((E \setminus F) \cup (E \cap F)) \\
        &= \mu(E \setminus F) + \mu(E \cap F) \\
        &\leq \mu(E \setminus F \cup F \setminus F) + \mu(E \cap F) \\
        &= \mu(E \triangle F) + \mu(E \cap F) \\
        &= \mu(E \cap F).
    \end{align*}
    De la misma forma, $\mu(F) = \mu(E \cap F) = \mu(E)$.

    \item Primero notemos que implícitamente asumimos que $E \triangle F \in \M$. Ahora, como $E \setminus F \cup F \setminus E = E \triangle F$ y $\mu(E \triangle F) = 0$, tenemos que $E \setminus F, F \setminus E \in \M$. Notemos que 
    \begin{align*}
        &(F \setminus E) \cup \cor{E \setminus (E \setminus F)} \\
        &= (F \cap E^c) \cup \cor{E \cap (E \cap F^c)^c} \\
        &= (F \cap E^c) \cup \cor{E \cap (E^c \cup F)} \\
        &= (F \cap E^c) \cup (E \cap F) \\
        &= F
    \end{align*}
    En particular tenemos 
    \[ F = \underbrace{(F \cap E^c)}_{\in\M} \cup \underbrace{(\underbrace{E}_{\in\M} \cap \underbrace{(\underbrace{E \cap F^c}_{\in\M})^c}_{\in\M})}_{\in\M} \in \M \]

    \item Notemos que $\mu(E \triangle E) = \mu(\varnothing) = 0 \implies E \sim E$. Por lo tanto es reflexiva. Notemos también que 
    \[ E \sim F \implies 0 = \mu(E \triangle F) = \mu(F \triangle E) \implies F \sim E. \]
    Luego para transitividad asumimos la parte 4, por lo que tenemos 
    \begin{align*}
        E \sim F \sim G
        &\implies \rho(E,F) = \rho(F,G) = 0  \\
        &\implies 0 \leq \rho(E,G) \leq \rho(E,F) + \rho(F,G) = 0 \\
        &\implies E \sim G.
    \end{align*}
    
    \item Notemos que 
    \begin{align*}
        \rho(E,F) &= \mu(E \triangle F) \\
        &= \mu(E \cap F^c) + \mu(E^c \cap F) \\
        &= \mu(E \cap F^c \cap G) + \mu(E \cap F^c \cap G^c) + \mu(E^c \cap F \cap G) + \mu(E^c \cap F \cap G^c) \\
        &\leq \mu(E^c \cap G) + \mu(E \cap G^c) + \mu(F^c \cap G) + \mu(F \cap G^c) \\
        &= \mu(E \triangle G) + \mu(F \triangle G) \\
        &= \rho(E,G) + \rho(F,G).
    \end{align*}
  \end{enumerate}
\end{proof}

\newpage
\subsection{Problema 4}

Sea $\P(\Z^+)$ la \( \sigma \)-álgebra en $\Z^+$ de todos los subconjuntos de $\Z^+$.

\begin{enumerate}
  \item Suponga que $\mu$ es una medida en $(\Z^+,\P(\Z^+))$. Muestre que existe una sucesión $(w_n)_n$ en $[0,\oo]$ tal que 
  \[ \mu(E) = \sum_{k \in E} w_k \]
  para cada $E \subseteq \Z^+$.
  \item De un ejemplo de alguna medida $\mu$ en $(\Z^+,\P(\Z^+))$ tal que 
  \[ \{\mu(E) \colon E \subseteq \Z^+\} = [0,1]. \]
\end{enumerate}


\newpage
\subsection{Problema 5}

Muestre que $B \subseteq \R$ es Lebesgue--medible si y solo si 
\[ \lam^*(I) = \lam^*(I \cap B) + \lam^*(I \cap B^c) \]
para todo intervalo abierto acotado $I \subseteq \R$.

\begin{proof}
  \ \begin{itemize}
    \item \boxed{\impliedby} Sea $B \subseteq \R$. Queremos 
    \[ \lam^*(B \cap A) + \lam^*(B \cap A^c) \leq \lam^*(B). \]
    Sea $\e > 0$. Sea $\{(a_n,b_n)\}_{n\in\N}$ subrimiento de $B$ tal que $\ds \sum_{n \in \N} (b_n - a_n) \leq \lam^*(B) + \e$. 
    \begin{align*}
        \lam^*(B) + \e &\geq \sum_{n\geq0} (b_n-a_n) \\
        &= \sum_{n\geq0} \lam^*(a_n,b_n) \\
        &= \sum_{n\geq0} \lam^*((a_n,b_n) \cap A) + \sum_{n\geq0} \lam^*((a_n,b_n) \cap A^c) \\
        &= \lam^*\br{\bigcup_{n\geq0} ((a_n,b_n) \cap A)} + \lam^*\br{\bigcup_{n\geq0} ((a_n,b_n) \cap A)} \\ 
        &= \lam^*\br{A \cap \bigcup_n (a_n,b_n)} + \lam^*\br{A^c \cap \bigcup_n (a_n,b_n)} \\
        &\geq \lam^*(B \cap A) + \lam^*(B \cap A^c)
    \end{align*}
    Esto para todo $\e>0$. Concluimos que $\lam^*(B \cap A) + \lam^*(B \cap A^c) \leq \lam^*(B)$.
  \end{itemize}
\end{proof}


\end{document}